\section{前向误差、后向误差与稳定算法}
\label{sec:08-Numerical-Analysis/01-Floating-Point-and-Error/04}

你在 review 一段解线性方程组的代码。测试是这么写的：算出 \(\hat x\) 之后，检查 \(\norm{A\hat x-b}\) 是不是小于 \(10^{-10}\)。这个测试从写下来那天起就没红过。

它确实在检查一件有意义的事，只是那件事不是``答案对不对''。残差小到机器精度的 \(\hat x\)，在一个病态系统上可以一位有效数字都不剩；而这个测试照样通过。

上一节把问题的敏感程度装进了 \(\kappa\)。现在轮到另一半：在给定精度下，你这套运算次序有没有制造出超出输入尺度的额外扰动。

\subsection{换一个问法：这个答案是谁的精确解}

设要算的是 \(y=f(x)\)，程序给出 \(\hat y\)。最直接的评价是前向误差 \(\norm{\hat y-f(x)}\)——算出来的答案离真答案多远。麻烦在于它几乎无法计算：要知道真答案，问题就已经解完了。

后向误差换了个问法：能不能找到一个邻近的输入 \(x+\Delta x\)，使得

\[
\hat y=f(x+\Delta x)
\]

精确成立？若所需的 \(\Delta x\) 很小，那么程序给出的东西是一个轻微扰动过的问题的精确解。

\begin{center}
\begin{tikzpicture}[x=1cm,y=1cm,font=\small,>=stealth]
  \node[font=\footnotesize,black!70] at (1.8,3.4) {输入空间};
  \node[font=\footnotesize,black!70] at (7.0,3.4) {输出空间};
  \fill[black!85] (1.8,2.5) circle (2.4pt);
  \fill[black!45] (1.8,0.9) circle (2.4pt);
  \fill[black!85] (7.0,2.7) circle (2.4pt);
  \fill[black!45] (7.0,0.7) circle (2.4pt);
  \node[above,font=\footnotesize] at (1.75,2.6) {\(x\)};
  \node[below,font=\footnotesize] at (1.85,0.8) {\(x+\Delta x\)};
  \node[right,font=\footnotesize] at (7.1,2.72) {\(f(x)\)};
  \node[right,font=\footnotesize] at (7.1,0.68) {\(\hat y\)};
  \draw[<->,red!70] (1.25,2.45) -- (1.25,0.95);
  \node[left,font=\footnotesize,red!75] at (1.2,1.7) {后向误差};
  \draw[<->,red!70] (7.9,2.7) -- (7.9,0.7);
  \node[right,font=\footnotesize,red!75] at (7.95,1.7) {前向误差};
  \draw[->,black!70] (2.0,2.52) -- (6.8,2.68);
  \node[above,font=\footnotesize,black!70] at (4.4,2.66) {精确的 \(f\)};
  \draw[->,black!70] (2.0,0.88) -- (6.8,0.72);
  \node[below,font=\footnotesize,black!70] at (4.4,0.74) {精确的 \(f\)};
  \draw[->,dashed,blue!70] (2.0,2.42) -- (6.85,0.75);
  \node[font=\footnotesize,blue!75] at (4.7,1.9) {算法（每一步都在舍入）};
\end{tikzpicture}

\emph{同一个 \(\hat y\)，两种读法。往右看，它离 \(f(x)\) 有多远；往下看，它是谁的精确解。右边那段距离难算，左边那段常常可以直接估出来。}
\end{center}

后向观点更贴合浮点计算的现实。输入本身多半已经由测量和舍入产生，要求程序忠于一个你从未真正拥有过的无限精确的数，这个要求没有意义；合理的要求是程序别再添乱。若算法在任何输入上的相对后向误差都是 \(O(u)\) 量级，就称它后向稳定。

最小的例子是一次浮点加法。上一节的模型给出 \(\operatorname{fl}(a+b)=(a+b)(1+\delta)\)，把右边的 \(1+\delta\) 分摊到两个加数上，它等于 \(a(1+\delta)+b(1+\delta)\)。算出来的和，正好是两个各自被相对扰动了 \(\abs{\delta}\le u\) 的输入的精确和。一次加法后向稳定，而且后向误差刚好卡在 \(u\)——这是任何算法能指望的最好情况。

\subsection{两笔账要相乘}

把两半接起来。若算法的相对后向误差是 \(O(u)\)，而问题在该点的相对条件数是 \(\kappa\)，那么图里那个邻近输入被 \(f\) 送到输出端时，位移被放大 \(\kappa\) 倍，于是一阶上

\[
\text{相对前向误差}\;\lesssim\;\kappa\times\text{相对后向误差}\;=\;O(\kappa u).
\]

条件数评价问题，稳定性评价算法，最终精度是两者的乘积。这句话展开就是一张四格表，它是这一章最该记住的东西。

\begin{center}
\begin{tikzpicture}[x=1cm,y=1cm,font=\small]
  \draw[fill=green!14,draw=black!45] (0,2.6) rectangle (4.4,5.2);
  \draw[fill=blue!12,draw=black!45] (4.4,2.6) rectangle (8.8,5.2);
  \draw[fill=orange!20,draw=black!45] (0,0) rectangle (4.4,2.6);
  \draw[fill=red!14,draw=black!45] (4.4,0) rectangle (8.8,2.6);
  \node[font=\footnotesize,align=center] at (2.2,3.9)
    {前向误差 \(\approx u\)\\满精度拿到手\\\emph{正常情形}};
  \node[font=\footnotesize,align=center] at (6.6,3.9)
    {残差小到 \(u\)，正确位数却很少\\误差 \(\approx\kappa u\)，已经到极限\\\emph{要动的是问题，不是代码}};
  \node[font=\footnotesize,align=center] at (2.2,1.3)
    {误差远大于 \(u\)\\全部由算法造成\\\emph{唯一能靠换算法修好的一格}};
  \node[font=\footnotesize,align=center] at (6.6,1.3)
    {两重放大叠在一起\\结果通常没有意义\\\emph{先修算法，再看还剩多少}};
  \node[font=\footnotesize] at (2.2,5.5) {问题良态（\(\kappa\) 小）};
  \node[font=\footnotesize] at (6.6,5.5) {问题病态（\(\kappa\) 大）};
  \node[rotate=90,font=\footnotesize] at (-0.35,3.9) {算法后向稳定};
  \node[rotate=90,font=\footnotesize] at (-0.35,1.3) {算法不稳定};
\end{tikzpicture}

\emph{横轴是问题的属性，纵轴是算法的属性，两者互不干涉。诊断的第一步是判断自己落在哪一格——左下角那格能靠调试解决，右上角那格无论怎么调都解决不了。}
\end{center}

四格里最容易被误读的是右上角。库函数报告残差 \(3\times10^{-16}\)，同时给出的解只有三位正确数字，这不矛盾：算法已经做到了极限，剩下的损失记在问题头上。而左下角是唯一值得重构代码的一格——良态问题上出现大误差，责任无处可推，要么算法不稳定，要么实现有错。

\subsection{残差是可以算出来的那一半}

线性系统把这套语言变得格外具体。对近似解 \(\hat x\)，残差 \(r=b-A\hat x\) 满足 \(A\hat x=b-r\)：\(\hat x\) 是右端被改成 \(b-r\) 的那个系统的精确解。相对残差 \(\norm{r}/(\norm{A}\norm{\hat x})\) 因此是一个可以直接算出来的后向误差指标。而前向误差是

\[
x-\hat x=A^{-1}r,
\]

要从残差推到它，必须乘上 \(A\) 的条件数。这正是上一节``小残差推不出小误差''在后向误差语言下的重述。

回到开头那个测试。它检查的是后向误差，这件事值得保留——它能抓住左下角那格，也就是实现层面的错误。但它证明不了答案接近真解。想让它有那个含义，得同时把 \(\kappa(A)\) 的估计报出来，让阅读测试结果的人自己乘一次。

至于稳定算法从哪儿来，规律相当集中：优先用正交变换，因为 \(\kappa_2\) 等于 1 的映射对误差既不放大也不压缩；选主元，别让接近零的数当除数；按量级缩放，把中间结果拉回不会溢出的区间；实在不行就换一个代数上等价的表达式，绕开会发生抵消的那一步——上一节二次方程先算大根、再用 \(x_1x_2=c/a\) 接出小根，用的就是最后这一招。

这些结论后面通常跟着几个默认条件，逐条对一遍能防住不同的失效。误差界里几乎总带着维数因子，形如 \(nu\) 或 \(\sqrt{n}\,u\)，\(n\) 上万时它不再可以忽略；Gauss 消去的界里另有一个增长因子，部分主元只把它压到经验上的 \(O(1)\)，最坏情形仍是 \(2^{n-1}\)，这是 Householder QR 在稳定性上更让人放心的理由，细节见\hyperref[sec:08-Numerical-Analysis/05-Numerical-Linear-Algebra/02]{``QR 与最小二乘的稳定解法''}；所有界都默认中间结果没有溢出、也没有掉进次正规区，而 \(\sqrt{a^{2}+b^{2}}\) 这种写法恰恰会先溢出，按 \(m=\max(\abs{a},\abs{b})\) 缩放成 \(m\sqrt{(a/m)^{2}+(b/m)^{2}}\) 就是为了防住这一条，\(m=0\) 时直接返回零以免出现 \(0/0\)；最后，一阶分析丢掉的高阶项要求 \(nu\ll1\)，否则整套估计本身就不成立。

\subsection{稳定这个词至少有四种用法}

数值分析里的``稳定''分属不同语境，混用会得出错误结论。浮点算法的前向或后向稳定说的是有限精度下的误差传播；常微分方程离散格式的稳定说的是扰动随时间步进是被压缩还是被放大；迭代法的稳定常常指是否收敛到不动点；统计里的稳健说的是分布偏移下估计量还靠不靠得住。

它们有联系，定义不能互换。一个随步长 \(h\to0\) 收敛的差分公式，在固定精度下把 \(h\) 取得过小反而会被舍入吃掉，最优步长落在截断与舍入两条曲线的交点上，这条权衡在\hyperref[sec:08-Numerical-Analysis/04-Numerical-Differentiation-and-Integration/01]{``数值微分为何不能无限缩小步长''}里算得很细。反过来，一个后向稳定的直接求解器压根不涉及收敛率——它没有迭代可言。

同样要防住的是拿样例下结论。几个良态例子算对，证明不了算法稳定，因为右上角和左上角在良态输入上表现一样；一个病态例子算错，也证明不了算法不稳定，那可能只是 \(\kappa u\) 本来就大于 1。要判断稳定性，得在良态问题上比对高精度参考，看误差是否停在 \(O(u)\)。

\subsection{两把尺子，量后面所有的算法}

这一章的四篇可以收成一句话带走：算得准不准，由问题的条件数与算法的后向误差相乘决定，而这两个因子必须分开去查。查错了因子，力气就白费——在病态问题上反复换求解器，和在良态问题上一味加精度，是同一种徒劳的两个方向。

后面章节里出现的每一个算法，都可以拿这两把尺子量一遍：它把误差压在哪个因子上，又对另一个因子无能为力。\hyperref[ch:065]{``数值线性代数：结构、稳定性与规模决定算法''}把这套量法用到了底——那里挑选分解方式的理由，几乎全部可以翻译成本节的四格表。
